% !TeX encoding = UTF-8
\documentclass[]{../../../../tex/classes/styledArticle}
\usepackage{../../../../tex/packages/hebrewSupport}
\usepackage{../../../../tex/packages/mathShortcuts}
\usepackage{../../../../tex/packages/theoremsSupport}

\usetikzlibrary{graphs}

\usepackage{listings}
\lstset{
	language=Python,
	numbers=left,
	breaklines,
	basicstyle=\ttfamily,
	commentstyle=\color{gray},
	prebreak=\textrightarrow,
}

\newcommand\squig{\rightsquigarrow}

\author{שחר פרץ}
\title{אלגוריתמים $\sim$ \textit{הרצאה 5}}
\begin{document}
	\maketitle
	\textbf{מרצה: }עמית ווינשטיין
	
	לאור כושר ההשרדות הנמוך שלנו בשיעור הקודם (שיעור הרדמות גבוה במיוחד בסוף השיעור), בשיעור הזה נסיים בשבע, יהיו יותר הפסקות, ונתמקד פחות בפורמליות של ההוכחות. 
	
	\section*{מסלולים קצרים ביותר}
	השיעור נדבר על מסלולים קצרים ביותר מנקודה אחת (באנגלית – single source shortest paths, בקיצור SSSP). הבעיה: 
	
	בהינתן גרף $G = (V, E)$ (מכוון ולא מכוון), ומשקלים $w \co E \to \R$ (מאפשרים משקלים שליליים). משקל של מסלול יוגדר להיות סכום משקליל הקשתות לאורכו. נתעניין במסלולים קלים ביותר (מק''בים). מסלול הוא לא בהכרח פשוט. 
	
	נבחין בבעיה – אם יש \textit{מעגל שלילי}, כלומר מעגל שמשקלו שלילי, נוכל לעבור בו כמה פעמים שנרצה מה שיאפשר למק''ב להיות במשקל $-\infty$. מכאן נוציא כמה נקודות חשובות: 
	\begin{itemize}
		\item נרצה שהאלגוריתמים שלנו יזהו ויעצרו אם יש מעגל שלילי. 
		\item בהינתן שאין מעגל שלילי, כל המק''בים יהיו מסלולים פשוטים (בה''כ – למעשה, מדויק יותר להגיד שיש להם גרסה פשוטה, בגלל מעגלים שאינם שליליים אך משקלם $0$). 
		\item לא ניתן להוסיף משקל קבוע לכל המשקלים שכן מסלולים שונים הם באורכים שונים. 
		\item אם $w$ היא הפונקציה הקבועה $1$, משקל המסלול=אורך המסלול פחות אחת. לכן BFS פותר את הבעיה בזמן לינארי. 
		\item סימון $\dg(s, v)$: משקל מהסלול קל ביותר מ־$s$ ל־$v$. במקרי קצה, $\dg(s, v) = \infty$ אם $v$ לא נגיש מ־$s$, ו־$\dg(s, v) = -\infty$ אם יש מעגל שלילי ''בדרך`` מ־$s$ ל־$v$. 
		\item בגרפים לא מכוונים, משקלים שליליים משמעותם זהה למשמעותם של מעגלים שליליים בגרפים מכוונים. 
	\end{itemize}
	
	ישנן תכונות נחמדות לסימון של $\dg(s, v)$: 
	\begin{itemize}
		\item א''ש המשולש: $\dg(s, v) \le \dg(s, u) + w(u, v)$. 
		\item תת־מסלול של מק''ב הוא מק''ב בפני עצמו. 
		\item אם $(u, v)$ קשת אחרונה במק''ב $s \squig v$, אזי $\dg(s, v) = \dg(s, u) + w(u, v)$. 
	\end{itemize}
	
	להלן הרעיון הכללי, ולאחר מכן נראה כל מיני אלגוריתמים ש''מממשים`` את הרעיון הזה: 
	
	
	\begin{itemize}
		\item נחזיק לכל קודקוד שני ערכים: 
		\begin{itemize}
			\item $\dd[v]$ – המשקל של המק''ב שמצאנו \textit{עד כה} מ־$s$ ל־$v$. 
			\item $\pi[v]$ – הקודקוד הקודם של $v$ במסלול, שמשקלו $\dd[v]$. 
		\end{itemize}
		\item פעולה בודדת \en{\texttt{Relax($u, v$)}} \ עבור קשת $(u, v)$ שבודקת האם אפשר לשפר מק''ב: 
		\sen\begin{lstlisting}
if d[u] + w(u, v) < d[u]:
	d[v] = d[u] + w(u, v)
	pi[v] = u\end{lstlisting}\she
		\item הרעיון: נתאחל את $\dd, \pi$ ונבצע \texttt{Relax} כמה שאפשר. 
	\end{itemize}
	ואז, האלגוריתמים שלנו יהיו שונים בצורה ובסדר בו הם עושים \texttt{Relax}־ים. 
	נבחין בכמה אבחנות בנוגע לגישה הזו: 
	\begin{itemize}
		\item בבירור $\dd[v] \ge \dg(s, v)$ לאורך כל האלגוריתמים. 
		\item ברגע ש־$\forall v \in V\co \dd[v] = \dg(s, v)$ ברור שלא ניתן לבצע יותר \texttt{Relax} שישפר משהו שכן זה בהכרח יקטין $\dd[v]$ כלשהו בסתירה לנקודה הקודמת. 
		\item בדומה ל־BFS, מצביעי $\pi$ במידה והמרחקים חושבו ייצגו מק''ב בהיפוך כיוון. 
	\end{itemize}
	
	נתחיל במקרה של גרף מכוון חסר מעגלים. 
	\begin{itemize}
		\item \textbf{אלגוריתם: }יש לגרף מיון טופולוגי. ניכר שקודקודים לפני $s$ אינם מעניינים (אי אפשר להגיע אליהם). עתה האלגוריתם פשוט: נעבור על הקודקודים לפי סדר המיון הטופולוגי, ולכל קודקוד נבצע \texttt{Relax} על כל הקשתות שלו. בסוף הריצה $\dd[v] = \dg(s, v)$ לכל $v \in V$. 
		\item \textbf{סיבוכיות: }$\oc(\sof V + \sof E)$
		\item \textbf{נכונות: }באינדוקציה על הסדר במיון הטופולוגי. נראה שכאשר מגיעים לקודקוד $i$, מתקיים $\dd[v_i] = \dg(s, v_i)$. נניח בה''כ ש־$s = v_0$ (לא משנה כי לפני $s$ מתקיים $\dd[v] = \infty$). עבור $i = 0$ הטענה מתקיימת ($\dd[s] = 0 = \dg(s, s)$). נניח נכונות עד $i - 1$ ונוכיח עד $i$. יהא $v_k$ קודקוד שלפני $v_i$ במסלול קל ביותר $s \squig v_i$. בהכרח $k < i$, ולכן $\dd[k] = \dg(s, k)$ בזמן הטיפול ב־$v_k$. בזמן ביצוע \en{\texttt{Relax($v_k$, $v_i$)}}\  נקבל מה.א.: 
		\[ \dd[v_i] \le \dd[v_k] + w(v_k, v_i) = \dg(s, v_k) + w(v_k, v_j) = \dg(s, v_i) \]
		כלומר $\dd[v_] = \dg(s, v_i)$. 
	\end{itemize}
	
	עתה נטפל במקרה הכללי. 
	\begin{itemize}
		\item \textbf{האסטרטגיה של Ford: }נבצע \texttt{Relax} עד שלא ניתן יותר. 
		\item מה יקרה אם יש מעגל שלילי? החרא הזה לא יעצור. 
		\item \textbf{שאלה: }כמה \texttt{Relax}־ים משפרים \textit{לכל היותר} נוכל לבצע? [איור של גרף אציקלי שבמקרה הגרוע נעשה בו $2^{\sof E}$ רילקסים. בגרף שבנה המורה בין כל זוג איברים במיון הטופולוגי יש קשת במשקל $0$ וקשת במשקל $2^{i}$]
		\item בגרף אציקלי, המסלולים הפשוטים הם באורך $\sof V - 1$. 
		\item \textbf{האסטרטגיה של Bellman-Ford: }במקום לעשות \texttt{Relax}־ים בלי שום סדר בינהם, ניקח את הקשתות בסדר כלשהו, ונעשה $\sof V - 1$ פעמים רצף \texttt{Relax}־ים על הקשתות לפי הסדר הזה. לאחר מכאן נבצע איטרציה אחרונה של \texttt{Relax}־ים, ואם משהו השתנה, נסיק שיש מעגל שלילי. 
		\item \textbf{האסטרטגיה של יואב עירוני: }אפשר לשלב את שני האלגוריתמים, לעשות מיון טופולוגי על גרף העל ולהשתמש באלגוריתם הראשון, ואז לטפל באמצעות בלמן־פורד בכל רכיב קשירות חזקה בנפרד. 
		\item \textbf{סיבוכיות Bellman-Ford: }$\oc(\sof E \cdot \sof V)$
		\item \textbf{נכונות: }אם אין מעגלים שליליים, אחרי $\sof V - 1$ איטרציות, נוכיח ש־$\dd[v] = \dg(s, v)$ לכל $v \in V$. נוכיח באינדוקציה על \textbf{אורך} (לא בהכרח משקל) מינימלי של מק''ב מ־$s$ לכל קודקוד $v$. עבור אורך $0$, $v = s$ וזה מתקיים. נוכיח באינדוקציה על האיטרציות שאחרי $i$ איטרציות כל קודקוד עם מק''ב באורך לכל היותר $i$ מקיים $\dd[v] = \dg(s, v)$. 
		
		אחרי $0$ איטרציות מתקיים $\dd[s] = 0 = \dg(s, )$ (אין מעגל שלילי). נניח עד $i -1$, נראה בעבור $i$. יהא $v$ קודקוד שאורך מינמלי במק''ב שלו הוא $i$. נניח $u$ קודקוד לפני $v$ במק''ב באורך זה. אורך מינמלי של מק''ב של $u$ הוא $i -1$ ולכן אחרי $i - 1$ איטרציות, מתקיים $\dd[u] = \dg(s, u)$. באיטרציה ה־$i$ נבצבע \texttt{Relax} על $u, v$ ונקבל $\dg(s, v) \le \dd[v] \le \dd[u] + w(u, v) = \dg(s, u) + w(u, v) = \dg(s, v)$ וסיימנו. 
		
		עתה נראה שהוא מזהה מעגלים שליליים ע''י בדיקת שינויים באיטרציה ה־$\sof V$. נניח שהאלגוריתם לא שיפר באיטרציה ה־$\sof V$. יהא $v_0 \dots v_k = v_0$ מעגל. נראה שבאיטרציה ה־$\sof V$ כל ה־\texttt{Relax}־ים לא שיפרו, כלומר $\forall i \co \dd[v_i] \le \dd[v_{i - 1}] + w(v_{i - 1}, v_i)$. נסכם את האי־שוויונים האלה: 
		\[ \sumnko \dd[v_i] \le \sum_{i = 0}^{k - 1}\dd[v_i] + \sumnko w(v_{i -1}, v_i) \implies 0 \le \sumnko w(v_{i - 1}, v_i) \]
		כלומר משקל המעגל אי־שלילי. 
		
		נניח שהיה שיפור באיטרציה ה־$\sof V$, נרצה לטעון שיש מעגל שלילי. אם לא היו מעגלים שליליים, כבר אחרי $\sof V - 1$ איטרציות התקיים $\dd[v] = \dg(s, v)$ לכל $v$. מצד שני, הלאוגיתם כן הצליח לבצע \texttt{Relax} כדי לשפר, נניח על קשת $(u, v)$. נשתמש בא''ש המשולש: 
		\[ \dd[u] + w(u, v) = \dg(s, u) + w(u, v) \ge \dg(s, v) = \dd[v] \]
		ולכן שום \texttt{Relax} אינו יכול לשפר. 
	\end{itemize}
	
	במקרה שבו אין משקלים שליליים, ניתן לבצע משהו יעיל יותר. 
	\begin{itemize}
		\item \textbf{האלגוריתם של Dijkstra: }האלגוריתם עובר על צמתים לפי סדר משקל עולה מ־$s$ ומבצעים \texttt{Relax} על הקשתות. נעזרים בתור עדיפויות על מנת לבחור את הצומת הבא בכל פעם. 
		\item \textbf{סיבוכיות: }נעשה $\oc(\sof E)$ של \texttt{Dec-Key}, נעשה $\oc(\sof V)$ הוספות לתור, ו־$\oc(\sof V)$ פעמים \texttt{Delete-Min}. הסיבוכיות היא תלוי מימוש. 
		\begin{itemize}
			\item בערימת פיבונאצי זה יצא $\oc(\sof V \log \sof V + \sof E)$. 
			\item בערימה רגילה $\oc(\cl{\sof V + \sof E}\log \sof V)$
			\item במערך $\oc(\sof V^2 + \sof E)$
		\end{itemize}
	\end{itemize}
	
	עתה נדבר על אלגוריתמים שמשפרים את הביצועים בעבור המקרה של חיפוש מק''ב בין $s$ ל־$t$ ספציפיים, אם כי לא אסימפטוטית. 
	\begin{itemize}
		\item \textbf{Bidirectional Dijkstra: }במקרה בו רוצים למצוא מסלול קל ביותר מ־$s$ ל־$t$ ספציפי. הרעיון: במקום לעשות רק Dijkstra מ־$s$, להפעיל בו זמנית גם מ־$s$ וגם מ־$t$, ולטעון שאם שתי הריצות נפגשו באמצע בקודקוד, הוא אינדיקציה לכך שמצאנו מסלול קצר ביותר (הוא לא בהכרח יהיה חלק ממנו, אך הוא מודיע לנו על כך שיש לנו את המידע הנגרש). 
		
		אינטואיטיבית, אם נחפש מ־$s$ ובמקביל מאחורה מ־$t$ נוכל לעיתים לחסוך בזמן החיפוש. בדוגמה של 2d grid זה יחסוך פי $2$ (מנוסחת שטח מעגל). בכל צעד נוציא מאחד התורים, לפי אצל מי מהם יש מפתח קטן ביותר. 
		
		הטענה העיקרית לגבי תנאי העצירה של זה: ברגע שיש קודקוד שיצא גם מ־$s$ וגם מ־$t$, בהכרח כבר גילינו מסלול קל ביות רמ־$s$ ל־$t$ (לאו דווקא כזה שעובר בקודקוד זה). האינטואציה להוכחה היא הפעלה של א''ש המשולש (בגרסה מורחבת, עם $\dg(u, t)$ במקום $w(u, t)$). 
		\item \textbf{אלגוריתם $\bm{\mathrm A^*}$: }אומנם מתחיל רק מ־$s$, אבל נותן עדיפות לדברים ש''יש סיכוי יותר טוב שיגיעו ל־$t$``. אי אפשר להפעיל אותו בכל גרף. נשתמש בפונקציית הערכה למרחק (שמתעלמת ממכושלים באמצע) מכל קודקוד $v$ ל־$t$ שמקיימת שתי תכונות: 
		\begin{itemize}
			\item הפונקציה חייבת להיות אופטימית, כלומר קטנה או שווה מהמרחק האמיתי. 
			\item צריכה לקיים את א''ש המשולש: $h(v) \le h(u) + w(u, v)$. 
		\end{itemize}
				במידה ו־$h = 0$, שקול ללהריץ Dijkstra. 
	\end{itemize}
	
	\textbf{שאלה: }כשאין משקלים שליליים, האם עץ מסלולים קלים ביותר מ־$s$ הוא בהכרח עץ פורש מינימלי. 
	
	תשובה: לא
	
	עתה ההרצאה תהפוך לתרגול. 
	
	\exe{נתונים $n$ סוגי מטבעות $c_1 \dots c_n$ וטבלת המרה $a_{ij}$ – כמה מטבעות מסוג $c_j$ נקבל בעבור מטבע מסוג $c_i$ (לדוגמה, נקבל 2.9 שקלים בעבור דולר אחד. ברשימות של המרצה/מתרגל עוד כתוב 3.6). נרצה לבדוק האם יש ארביטראג', כלומר, סדרת המרות שבסופה נשאר עם יותר כסף ממה שהתחלנו. }
	\textbf{הרעיון: }אנחנו מחפשים רצף מטבעות (הכפלות בשערים) $c_{i_1} \dots c_{i_k}$ כך שאם נתחיל ממטבע מסוג $c_{i_1}$ נסיים עם יותר ממנו בסוף, כלומר, נרצה ש־$a_{i_1, i_2} \cdot a_{i_2, i_3} \cdots a_{i_k, i_1} > 1$. נוציא $\log$ כך שנרצה למעשה שסכום הלוגים יהיה גדול מ־$0$. 
	
	נבנה גרף שמשקליו הם $-\log a_{i, j}$ עבור הקשת המכוונת $(i, j)$ (בין המטבע ה־$i$ למטבע ה־$j$) ונחפש מעגל שלילי. סיבוכיות: $\oc(n^{3})$. 
	
	\exe{נתונה מערכת א''שים מהצורה $x_i - x_j \le c_{ij}$ עבור משתנים $x_1 \dots x_n$. נרצה פתרון למערכת, אם קיים, או להגיד שאין כזה. }
	
	\textbf{הפתרון: }נבנה גרף עם $n + 1$ קודקודים, $x_1 \dots x_n$ וקדקוד $s$ שמחובר לכולם עם משקל $0$. נוסיף עבור אי־השוויון $x_i - x_j \le c_{ij}$ קשת $(x_j, x_i)$ במשקל $c_{ij}$. נמצא מרחקים מ־$s$ ע''י BF. אם יש מעגל שלילי, אין פתרון. אחרת, $\dg(s, x_i)$ הוא הצבה טובה עבור $x_i$. זאת כי: 
	\[ \dg(s, x_i) \le \dg(s, x_j) + \c_{ij} \iff x_i - x_j \le c_{ij} \]
	ומעגל שלילי אמ''מ קב' אש''ים שהסכומים שלהם גוררים סתירה. 
	
	\textbf{שאלה: }גרף מכוון ללא מעגלים שליליים. נרצה לבדוק אם יש בו מעגל ממשקל $0$. 
	
	\textbf{רעיון: }נוסיף קודקוד $s$ עם קשת במשקל $0$ לכל שאר הקודקודים. נריץ BF מ־$s$. נבנה גרף מק''בים מ־$s$, כלומר, נכלול רק קשתות $(u, v)$ עבורן $\dg(s, v) = \dg(s, u) + w(u, v)$. נחפש מעגל בגרף זה (למשל DFS) ונטען שקיים מעגל בגרף זה אמ''מ יש מעגל במשקל $0$ בגרף המקורי. זה הטריק שראינו מקודם של סכימה וחיסור סכום ה־$\dg(s, v)$־ים משני האגפים. 
	
	\textbf{שאלה: }נתון גרף $G$ מכוון  וכל המשקלים מהקבוצה $\{0, 1, 2\}$. רוצים למצוא מק''ב מ־$S$. 
	
	\textbf{פתרון: }אפשר להחזיר טור עדיפויות של מערך באורך $2\sof V$ במקום פיבונאצי. ואז סך פעולות \texttt{Del-Min} יעלו $\oc(\sof V)$ וסך פעולות \texttt{Dec-Key} יעלו $\oc(\sof E)$. קודקודים ישמשו כמפתח במערך,, אך פעם לא חוזרים אחורה. 
	
	\textbf{שאלה: }מצאו מק''בים באורך זוגי. 
	
	\textbf{רעיון: }כמו פעם קודמת אפשר להשתמש בגרף דו''צ מצחיק. 
	
	\textbf{שאלה: }איך מוצאים מק''בים כאשר יש משקל גם על צמתים? 
	
	\textbf{פתרון: }נחליף כל צומת בשני צמתים עם קשת עליה יש את המשקל של הצומת המקורית, ואז נפעיל את האלגוריתם הרגיל. 
	
	\ndoc
\end{document}